
\section{The XSum Dataset}
\label{sec:bbcdataset}

% \begin{table*}[ht]
%   \begin{center}{\footnotesize
%   \begin{tabular}{ l | l | c c | c c | r r}
%     \thickhline
%     \multirow{2}{*}{Datasets} & \multirow{2}{*}{\# docs (train/val/test)} & \multicolumn{2}{c|}{avg. document length} & \multicolumn{2}{c|}{avg. summary length} & \multicolumn{2}{c}{vocabulary size}\\
%     & & words & sentences & words & sentences & document & summary\\ \thickhline 
%     CNN & 90,266/1,220/1,093 & 760.50 & 33.98 & 45.70 & 3.59 & 343,516 & 89,051 \\
%     DailyMail & 196,961/12,148/10,397 & 653.33 & 29.33 & 54.65 & 3.86 & 563,663 & 179,966 \\
%     NY Times & 589,284/32,736/32,739 & 800.04 & 35.55 & 45.54 & 2.44 & 1,399,358 & 294,011 \\ 
%     Newsroom & 992,966/108,591/108,650 & 770.09 & 34.73 & 30.36 & 1.43 & 2,646,681 & 360,290\\ \hline
%     Newsroom-ext & 331,778/36,332/36,122 & 706.06 & 31.65 & 45.78 & 1.88 & 1,214,748 & 243,062 \\
%     Newsroom-mixed & 328,634/35,879/36,006 & 830.58 & 36.63 & 23.78 & 1.17 & 1,271,435 & 169,875 \\
%     Newsroom-abs & 332,554/36,380/36,522 & 774.17 & 35.92 & 21.49 & 1.25 & 1,385,205 & 157,939 \\ \hline
%     XSum & 204,045/11,332/11,334 & 431.07 & 19.77 & 23.26 & 1.00 & 399,147 & 81,092 \\ \thickhline
%   \end{tabular}}
%   \end{center}
%   \caption{We compare  datasets
%     with respect to overall corpus size, size of training, validation,
%     and test sets, average document (source) and summary (target)
%     length (in terms of words and sentences), and  vocabulary size on
%     both on source and target. For CNN and DailyMail,
%     we  used the original splits of
%     \protect\newcite{hermann-nips15} and followed
%     \protect\newcite{narayan-rank18} to preprocess them. For  NY
%     Times  \protect\cite{nytcorpus}, we  used the splits
%     and pre-processing steps of
%     \protect\newcite{paulus-socher-arxiv17}. For the Newsroom dataset, we use the splits
%     and pre-processing steps of \protect\newcite{newsroom-naacl18}. We present the full 
%     Newsroom dataset (Newsroom) along with its three subsets: mostly extractive (Newsroom-ext), 
%     mixed (Newsroom-mixed) and mostly abstractive (Newsroom-abs).
%     For the vocabulary, we
%     lower case tokens. } \label{table:bbc-size-comparison} 
% \end{table*}

\begin{table*}
  \begin{center}{\footnotesize
  \begin{tabular}{ l | l | c c | c c | r r}
    \thickhline
    \multirow{2}{*}{Datasets} & \multirow{2}{*}{\# docs (train/val/test)} & \multicolumn{2}{c|}{avg. document length} & \multicolumn{2}{c|}{avg. summary length} & \multicolumn{2}{c}{vocabulary size}\\
    & & words & sentences & words & sentences & document & summary\\ \thickhline 
    CNN & 90,266/1,220/1,093 & 760.50 & 33.98 & 45.70 & 3.59 & 343,516 & 89,051 \\
    DailyMail & 196,961/12,148/10,397 & 653.33 & 29.33 & 54.65 & 3.86 & 563,663 & 179,966 \\
    NY Times & 589,284/32,736/32,739 & 800.04 & 35.55 & 45.54 & 2.44 & 1,399,358 & 294,011 \\
    XSum & 204,045/11,332/11,334 & 431.07 & 19.77 & 23.26 & 1.00 & 399,147 & 81,092 \\ \thickhline
  \end{tabular}}
  \end{center}
  \caption{Comparison of summarization datasets
    with respect to overall corpus size, size of training, validation,
    and test set, average document (source) and summary (target)
    length (in terms of words and sentences), and  vocabulary size on
    both on source and target. For CNN and DailyMail,
    we  used the original splits of
    \protect\newcite{hermann-nips15} and followed
    \protect\newcite{narayan-rank18} to preprocess them. For  NY
    Times  \protect\cite{nytcorpus}, we  used the splits
    and pre-processing steps of
    \protect\newcite{paulus-socher-arxiv17}. For the vocabulary, we
    lowercase tokens.} \label{table:bbc-size-comparison} 
\end{table*}

% \begin{table*}
%   \begin{center}{\footnotesize
%   \begin{tabular}{ l | c c c c | c c c | c c c } %  l l l l | l l l | l l l}
%     \thickhline
%     \multirow{2}{*}{Datasets} & \multicolumn{4}{c|}{\% of novel n-grams in gold summary} & \multicolumn{3}{c|}{\textsc{lead}} & \multicolumn{3}{c}{\textsc{ext-oracle}} \\
%     & unigrams & bigrams & trigrams & 4-grams & R1 & R2 & RL & R1 & R2 & RL \\ \thickhline 
%     CNN & 16.75 & 54.33 & 72.42 & 80.37 & 29.15 & 11.13 & 25.95 & 50.38 & 28.55 & 46.58 \\
%     DailyMail & 17.03 & 53.78 & 72.14 & 80.28 & 40.68 & 18.36 & 37.25 & 55.12 & 30.55 & 51.24 \\
%     NY Times & 22.64 & 55.59 & 71.93 & 80.16 & 31.85 & 15.86 & 23.75 & 52.08 & 31.59 & 46.72 \\ 
%     Newsroom & 18.31 & 46.80 & 58.06 & 62.72 & 33.04 & 22.35 & 30.31 & 57.09 & 42.94 & 53.65 \\ \hline
%     Newsroom-ext & 2.65 & 7.25 & 10.25 & 12.42 & 55.87 & 50.60 & 54.76 & 89.63 & 87.20 & 89.32 \\
%     Newsroom-mixed & 13.78 & 48.37 & 67.15 & 77.11 & 27.95 & 13.87 & 23.97 & 51.98 & 34.04 & 46.96 \\ 
%     Newsroom-abs & \textbf{38.25} & \textbf{84.36} & \textbf{96.39} & 98.28 & \textbf{15.44} & 2.72 & 12.32 & 29.85 & \textbf{7.82} & 24.86\\ \hline
%     XSum & 35.76 & 83.45  & 95.50  & \textbf{98.49}  & 16.30 & \textbf{1.61}  & \textbf{11.95}  & \textbf{29.79}  & 8.81  & \textbf{22.65}  \\ \thickhline
%   \end{tabular}}
%   \end{center}
%   \caption{Corpus bias towards extractive methods in the CNN,
%     DailyMail, NY Times, Newsroom, and XSum datasets.  We show  the proportion of novel $n$-grams in gold summaries. We also report ROUGE scores  for
%     the  \textsc{lead} baseline and 
%     \textsc{ext-oracle}, the extractive oracle system.     All results
%     are computed on the test set.} \label{table:ngram-coverage-lead-oracle}
% \end{table*}

\begin{table*}
  \begin{center}{\footnotesize
  \begin{tabular}{ l | c c c c | c c c | c c c } %  l l l l | l l l | l l l}
    \thickhline
    \multirow{2}{*}{Datasets} & \multicolumn{4}{c|}{\% of novel n-grams in gold summary} & \multicolumn{3}{c|}{\textsc{lead}} & \multicolumn{3}{c}{\textsc{ext-oracle}} \\
    & unigrams & bigrams & trigrams & 4-grams & R1 & R2 & RL & R1 & R2 & RL \\ \thickhline 
    CNN & 16.75 & 54.33 & 72.42 & 80.37 & 29.15 & 11.13 & 25.95 & 50.38 & 28.55 & 46.58 \\
    DailyMail & 17.03 & 53.78 & 72.14 & 80.28 & 40.68 & 18.36 & 37.25 & 55.12 & 30.55 & 51.24 \\
    NY Times & 22.64 & 55.59 & 71.93 & 80.16 & 31.85 & 15.86 & 23.75 & 52.08 & 31.59 & 46.72 \\
    XSum & \textbf{35.76} & \textbf{83.45}  & \textbf{95.50}  & \textbf{98.49}  & \textbf{16.30} & \textbf{1.61}  & \textbf{11.95}  & \textbf{29.79}  & \textbf{8.81}  & \textbf{22.65}  \\ \thickhline
  \end{tabular}}
  \end{center}
  \caption{Corpus bias towards extractive methods in the CNN,
    DailyMail, NY Times, and XSum datasets.  We show  the proportion of novel $n$-grams in gold summaries. We also report ROUGE scores  for
    the  \textsc{lead} baseline and the extractive oracle system
    \textsc{ext-oracle}.     Results
    are computed on the test set.
    \label{table:ngram-coverage-lead-oracle}}
\end{table*}


%BBC articles often have
%their first sentence, tagged with an HTML class
%``story-body\_\_introduction,'' written to introduce the article.
%Our aim is to predict that sentence from the rest of the
%article.
%The first sentence in the BBC articles include important information from across the
%whole document, and as such generating it requires inference across the document.
%In addition, since we have one sentence at hand, generating it requires
%complex fusion, synthesis and paraphrasing. Figure \ref{fig:bbcex-1}
%shows an example article from our dataset exhibiting these operations.


Our extreme summarization dataset (which we call XSum) consists of BBC
articles and accompanying single sentence summaries. Specifically,
each article is prefaced with an introductory sentence (aka summary)
which is professionally written, typically by the author of the
article. The summary bears the HTML class
``story-body\_\_introduction,'' and can be easily identified and
extracted from the main text body (see Figure~\ref{fig:bbcex-1} for an
example summary-article pair).



% Discuss the example to promote the task. 

%\paragraph{Data Collection and Statistics.} 


We followed the methodology proposed in \newcite{hermann-nips15} to
create a large-scale dataset for extreme summarization. Specifically,
we collected~226,711 Wayback archived BBC articles ranging over almost
a decade (2010 to 2017) and covering a wide variety of domains
(e.g.,~News, Politics, Sports, Weather, Business, Technology, Science,
Health, Family, Education, Entertainment and Arts). Each article comes
with a unique identifier in its URL, which we used to randomly split
the dataset into training (90\%, 204,045), validation (5\%, 11,332),
and test (5\%, 11,334) set. Table~\ref{table:bbc-size-comparison}
compares XSum with the CNN, DailyMail, and NY Times % , and Newsroom
benchmarks. As can be seen, XSum contains a substantial number of
training instances, similar to DailyMail; documents and summaries in
XSum are shorter in relation to other datasets but the vocabulary size
is sufficiently large, comparable to CNN.

% grep 'web/2017' map-webarxive-bbcid.txt | wc -l 119262
% grep 'web/2016' map-webarxive-bbcid.txt | wc -l 70457
% grep 'web/2015' map-webarxive-bbcid.txt | wc -l 44361 
% grep 'web/2014' map-webarxive-bbcid.txt | wc -l 2935
% grep 'web/2013' map-webarxive-bbcid.txt | wc -l 1581
% grep 'web/2012' map-webarxive-bbcid.txt | wc -l 928
% grep 'web/2011' map-webarxive-bbcid.txt | wc -l 221
% grep 'web/2010' map-webarxive-bbcid.txt | wc -l 43

%\paragraph{Reduced Bias towards Extractive Methods.} 
Table~\ref{table:ngram-coverage-lead-oracle} provides empirical
analysis supporting our claim that XSum is less biased toward
extractive methods compared to other summarization datasets.  We
report the percentage of novel $n$-grams in the target gold summaries
that do not appear in their source documents.  There are 36\% novel
unigrams in the XSum reference summaries compared to 17\% in CNN,
17\%~in DailyMail, and 23\%~in NY Times. % , and 18\%~in Newsroom.
This indicates that XSum summaries are more abstractive. The
proportion of novel constructions grows for larger $n$-grams across
datasets, however, it is much steeper in XSum whose summaries exhibit
approximately 83\% novel bigrams, 96\% novel trigrams, and 98\% novel
4-grams (comparison datasets display around 47--55\%~new bigrams,
58--72\%~new trigrams, and 63--80\%~novel 4-grams).

We further evaluated two extractive methods on these datasets.
\textsc{lead} is often used as a strong lower bound for news
summarization \cite{nenkova-05} and creates a summary by selecting the
first few sentences or words in the document. We extracted the first
3~sentences for CNN documents and the first 4~sentences for DailyMail
\cite{narayan-rank18}. Following previous work
\cite{durrett-nyt-ext,paulus-socher-arxiv17}, we obtained
\textsc{lead} summaries based on the first 100 words for NY Times
documents. % For Newsroom, we extracted the first 2~sentences to form
% the \textsc{lead} summaries.
For XSum, we selected the first sentence
in the document (excluding the one-line summary) to generate the
\textsc{lead}. Our second method, \textsc{ext-oracle}, can be viewed
as an upper bound for extractive models
\cite{nallapati17,narayan-rank18}. It creates an oracle summary by
selecting the best possible set of sentences in the document that
gives the highest ROUGE \cite{rouge} with respect to the gold
summary. For XSum, we simply selected the single-best sentence in the
document as summary.

Table~\ref{table:ngram-coverage-lead-oracle} reports the performance
of the two extractive methods using ROUGE-1 (R1), ROUGE-2 (R2), and
ROUGE-L (RL) with the gold summaries as reference. The \textsc{lead}
baseline performs extremely well on CNN, DailyMail and NY
Times % and Newsroom
confirming that they are biased towards extractive methods.
\textsc{ext-oracle} further shows that improved sentence selection
would bring further performance gains to extractive
approaches. Abstractive systems trained on these datasets often have a
hard time beating the \textsc{lead}, let alone \textsc{ext-oracle}, or
display a low degree of novelty in their summaries
\cite{see-acl17,tanwan-acl17,paulus-socher-arxiv17,Pasunuru-multireward18,asli-multiagent18}.
Interestingly, \textsc{lead} and \textsc{ext-oracle} perform poorly on
XSum underlying the fact that it is less biased towards extractive
methods.


%Our results demonstrate that the XSum dataset encourages more fusion
%and paraphrasing to generate summaries. In fact, our novel abstractive
%summarization model accommodating these features (see Section
%\ref{sec:topicconvabssum}), outperforms both \textsc{lead} and
%\textsc{ext-oracle} on this dataset. As a result, we believe that the
%XSum dataset provides a better benchmark for training and testing
%abstractive summarization systems. 


In line with our findings, \newcite{newsroom-naacl18} have recently
reported similar extractive biases in existing datasets. They
constructed a new dataset called ``Newsroom'' which demonstrates a
high diversity of summarization styles. XSum is not diverse, it
focuses on a single news outlet (i.e.,~BBC) and a unifrom
summarization style (i.e.,~a single sentence). However, it is
sufficiently large for neural network training and we hope it will
spur further research towards the development of abstractive
summarization models.

% In line with our findings, \newcite{newsroom-naacl18} have recently
% reported similar extractive biases in existing datasets. Their newly
% constructed ``Newsroom'' dataset demonstrates a high diversity of
% summarization styles: mostly extractive, mixed and mostly
% abstractive. The abstractive subset of Newsroom (Newsroom-abs)
% demonstrates similar patterns to XSum in terms of the percentage of
% novel n-grams in gold summary and the performances of extractive
% methods, \textsc{lead} and \textsc{ext-oracle}. However, XSum differs
% from Newsroom in the way these summaries are collected. Newsroom
% summaries are extracted using the HTML meta-tag {\em description} from
% a wide range of news outlets, these descriptions are often used for
% article indexing and often not shown to the readers. % These
% % descriptions could be noisy and ungrammatical.
% In contrast, XSum focuses on a single news outlet (i.e.,~BBC) and a
% unifrom summarization style (i.e.,~a single sentence). XSum summaries
% are always shown to the readers along with the article.  XSum is
% sufficiently large for neural network training and we hope it will
% spur further research towards the development of abstractive
% summarization models.






%In this paper, however, we propose
%the XSum benchmark that encourages high levels of abstraction. In
%addition, we propose a novel abstractive summarization model (see the
%following section) more suitable for this task.

%%% Local Variables:
%%% mode: latex
%%% TeX-master: "extreme-sum"
%%% End:

